% ============================================================================
% © 2025 Mohamed EL AFRIT — IPSSI
% Licence CC BY-NC-SA 4.0
% https://creativecommons.org/licenses/by-nc-sa/4.0/deed.fr
%
% FICHE D'EXERCICES (énoncés) — Jour 1
% Données et représentation mathématique
% ============================================================================
% ============================================================================
% © 2025 Mohamed EL AFRIT — IPSSI
% Ce contenu est distribué sous licence Creative Commons BY-NC-SA 4.0
% https://creativecommons.org/licenses/by-nc-sa/4.0/deed.fr
% ============================================================================
%
% TEMPLATE COMMUN — Mathématiques appliquées à l'Intelligence Artificielle
% Ce fichier est le préambule partagé de TOUS les documents LaTeX du projet.
%
% Utilisation :
%   % ============================================================================
% © 2025 Mohamed EL AFRIT — IPSSI
% Ce contenu est distribué sous licence Creative Commons BY-NC-SA 4.0
% https://creativecommons.org/licenses/by-nc-sa/4.0/deed.fr
% ============================================================================
%
% TEMPLATE COMMUN — Mathématiques appliquées à l'Intelligence Artificielle
% Ce fichier est le préambule partagé de TOUS les documents LaTeX du projet.
%
% Utilisation :
%   % ============================================================================
% © 2025 Mohamed EL AFRIT — IPSSI
% Ce contenu est distribué sous licence Creative Commons BY-NC-SA 4.0
% https://creativecommons.org/licenses/by-nc-sa/4.0/deed.fr
% ============================================================================
%
% TEMPLATE COMMUN — Mathématiques appliquées à l'Intelligence Artificielle
% Ce fichier est le préambule partagé de TOUS les documents LaTeX du projet.
%
% Utilisation :
%   \input{template_ipssi}          % Importer ce préambule
%   \jourNumero{1}                  % Définir le numéro du jour (1, 2, 3 ou 4)
%   \jourTheme{Données et représentation mathématique}
%   \begin{document}
%   \pagegarde{Fiche de synthèse}{1}{Données et représentation mathématique}
%   ...
%   \end{document}
% ============================================================================

% --- Classe de document ---
\documentclass[11pt,a4paper]{article}

% ============================================================================
% ENCODAGE ET LANGUE
% ============================================================================
\usepackage[utf8]{inputenc}
\usepackage[T1]{fontenc}
\usepackage[french]{babel}

% ============================================================================
% MATHÉMATIQUES
% ============================================================================
\usepackage{amsmath,amssymb,amsfonts,mathtools}

% Commandes mathématiques personnalisées (conventions du cours)
\newcommand{\vx}{\mathbf{x}}          % vecteur x
\newcommand{\vw}{\mathbf{w}}          % vecteur w (poids)
\newcommand{\vy}{\mathbf{y}}          % vecteur y
\newcommand{\mX}{\mathbf{X}}          % matrice X
\newcommand{\mW}{\mathbf{W}}          % matrice W
\newcommand{\yhat}{\hat{y}}           % prédiction
\newcommand{\pscal}[2]{\langle #1, #2 \rangle} % produit scalaire
\newcommand{\gradJ}{\nabla J(\boldsymbol{\theta})} % gradient de J

% ============================================================================
% MISE EN PAGE
% ============================================================================
\usepackage[margin=2.5cm]{geometry}
\usepackage{setspace}
\setstretch{1.15}                      % interligne légèrement aéré

% ============================================================================
% COULEURS
% ============================================================================
\usepackage[dvipsnames,svgnames,x11names]{xcolor}

% Palette principale du projet
\definecolor{ipssiBleu}{HTML}{3B82F6}       % Définitions
\definecolor{ipssiVert}{HTML}{22C55E}       % À retenir
\definecolor{ipssiOrange}{HTML}{F97316}     % Attention / Piège
\definecolor{ipssiViolet}{HTML}{A855F7}     % Intuition
\definecolor{ipssiCyan}{HTML}{06B6D4}       % Exemples
\definecolor{ipssiGrisClair}{HTML}{F3F4F6}  % Fond code
\definecolor{ipssiGrisFonce}{HTML}{1F2937}  % Texte code
\definecolor{ipssiFondPage}{HTML}{FAFAFA}   % Fond léger page de garde
\definecolor{ipssiRouge}{HTML}{E74C3C}      % Classe 0 (salaire ≤ 45k€)
\definecolor{ipssiVertData}{HTML}{2ECC71}   % Classe 1 (salaire > 45k€)
\definecolor{ipssiBleuDroite}{HTML}{3498DB} % Droite de régression

% ============================================================================
% ENCADRÉS PÉDAGOGIQUES (tcolorbox)
% ============================================================================
\usepackage[most]{tcolorbox}

% Style de base partagé
\tcbset{
    ipssibase/.style={
        enhanced,
        breakable,
        left=4mm,
        right=4mm,
        top=3mm,
        bottom=3mm,
        boxrule=0pt,
        frame hidden,
        borderline west={3pt}{0pt}{#1},
        colback=#1!5,
        colframe=#1,
        fonttitle=\bfseries\sffamily,
        coltitle=#1!80!black,
        attach boxed title to top left={yshift=-2mm, xshift=4mm},
        boxed title style={
            colback=#1!15,
            colframe=#1,
            boxrule=0.5pt,
            arc=2pt,
        },
        arc=2pt,
        outer arc=2pt,
    }
}

% Icônes TikZ pour les encadrés (compatible pdflatex, pas d'emoji Unicode)
\newcommand{\icoDefinition}{%
    \tikz[baseline=-0.3em]{\draw[ipssiBleu, thick] (0,0) -- (0.3,0) -- (0.3,0.35) -- (0,0.35) -- cycle;
    \draw[ipssiBleu, thick] (0.05,0.1) -- (0.25,0.1); \draw[ipssiBleu, thick] (0.05,0.2) -- (0.2,0.2);}%
}
\newcommand{\icoRetenir}{%
    \tikz[baseline=-0.2em]{\draw[ipssiVert, thick, fill=ipssiVert!20] (0.15,0) -- (0.3,0.3) -- (0,0.3) -- cycle;}%
}
\newcommand{\icoAttention}{%
    \tikz[baseline=-0.2em]{\draw[ipssiOrange, thick] (0.15,0.35) -- (0,0) -- (0.3,0) -- cycle;
    \node[ipssiOrange, font=\tiny\bfseries] at (0.15,0.12) {!};}%
}
\newcommand{\icoIntuition}{%
    \tikz[baseline=-0.2em]{\draw[ipssiViolet, thick, fill=ipssiViolet!15]
    (0.15,0.35) -- (0.08,0.2) -- (0.08,0.15) -- (0.22,0.15) -- (0.22,0.2) -- cycle;
    \draw[ipssiViolet, thick] (0.11,0.12) -- (0.11,0.05); \draw[ipssiViolet, thick] (0.19,0.12) -- (0.19,0.05);}%
}
\newcommand{\icoExemple}{%
    \tikz[baseline=-0.2em]{\draw[ipssiCyan, thick] (0,0) rectangle (0.3,0.3);
    \draw[ipssiCyan, thick] (0,0.15) -- (0.3,0.15); \draw[ipssiCyan, thick] (0.15,0) -- (0.15,0.3);}%
}
\newcommand{\icoPython}{%
    \tikz[baseline=-0.2em]{\node[font=\small\bfseries\ttfamily, ipssiVert!70!black] at (0,0) {\textgreater\textgreater};}%
}

% Définition (bleu)
\newtcolorbox{definitionbox}[1][D\'{e}finition]{
    ipssibase=ipssiBleu,
    title={\icoDefinition\; #1},
}

% À retenir (vert)
\newtcolorbox{retenirbox}[1][\`{A} retenir]{
    ipssibase=ipssiVert,
    title={\icoRetenir\; #1},
}

% Attention / Piège (orange)
\newtcolorbox{attentionbox}[1][Attention]{
    ipssibase=ipssiOrange,
    title={\icoAttention\; #1},
}

% Intuition (violet)
\newtcolorbox{intuitionbox}[1][Intuition]{
    ipssibase=ipssiViolet,
    title={\icoIntuition\; #1},
}

% Exemple (cyan)
\newtcolorbox{exemplebox}[1][Exemple]{
    ipssibase=ipssiCyan,
    title={\icoExemple\; #1},
}

% Code Python (style terminal)
\newtcolorbox{pythonbox}[1][Code Python]{
    enhanced,
    breakable,
    left=4mm, right=4mm, top=3mm, bottom=3mm,
    boxrule=0.5pt,
    colback=ipssiGrisClair,
    colframe=ipssiGrisFonce!40,
    fonttitle=\bfseries\ttfamily\small,
    coltitle=ipssiGrisFonce,
    title={\icoPython\; #1},
    attach boxed title to top left={yshift=-2mm, xshift=4mm},
    boxed title style={colback=ipssiGrisClair, colframe=ipssiGrisFonce!40, boxrule=0.5pt, arc=2pt},
    arc=2pt, outer arc=2pt,
}

% ============================================================================
% CODE PYTHON (listings)
% ============================================================================
\usepackage{listings}

\lstdefinestyle{python}{
    language=Python,
    basicstyle=\ttfamily\small,
    keywordstyle=\color{ipssiBleu}\bfseries,
    stringstyle=\color{ipssiVert!70!black},
    commentstyle=\color{gray}\itshape,
    numberstyle=\tiny\color{gray},
    numbers=left,
    numbersep=8pt,
    backgroundcolor=\color{ipssiGrisClair},
    frame=none,
    rulecolor=\color{ipssiGrisFonce!20},
    breaklines=true,
    breakatwhitespace=true,
    tabsize=4,
    showstringspaces=false,
    captionpos=b,
    aboveskip=8pt,
    belowskip=8pt,
    xleftmargin=10pt,
    framexleftmargin=10pt,
    morekeywords={as, with, True, False, None, self, np, plt, import, from},
    literate=
        {é}{{\'e}}1 {è}{{\`e}}1 {ê}{{\^e}}1 {ë}{{\"e}}1
        {à}{{\`a}}1 {â}{{\^a}}1
        {ù}{{\`u}}1 {û}{{\^u}}1
        {ô}{{\^o}}1
        {î}{{\^i}}1 {ï}{{\"i}}1
        {ç}{{\c{c}}}1
        {É}{{\'E}}1 {È}{{\`E}}1
        {→}{{$\rightarrow$}}1
        {≈}{{$\approx$}}1
        {≤}{{$\leq$}}1
        {≥}{{$\geq$}}1
        {★}{{$\bigstar$}}1
        {ŷ}{{$\hat{y}$}}1
        {·}{{$\cdot$}}1,
}

\lstset{style=python}

% Style pour le code inline
\newcommand{\pyinline}[1]{\lstinline[style=python,basicstyle=\ttfamily\small]{#1}}

% ============================================================================
% GRAPHIQUES
% ============================================================================
\usepackage{tikz}
\usepackage{pgfplots}
\pgfplotsset{compat=1.18}
\usetikzlibrary{arrows.meta, calc, positioning, shapes.geometric, decorations.pathreplacing}

% ============================================================================
% TABLEAUX
% ============================================================================
\usepackage{booktabs}
\usepackage{tabularx}
\usepackage{multirow}

% ============================================================================
% LISTES
% ============================================================================
\usepackage{enumitem}
\setlist[itemize]{itemsep=2pt, parsep=0pt, topsep=4pt}
\setlist[enumerate]{itemsep=2pt, parsep=0pt, topsep=4pt}

% ============================================================================
% HYPERLIENS
% ============================================================================
\usepackage[
    colorlinks=true,
    linkcolor=ipssiBleu!80!black,
    urlcolor=ipssiBleu!80!black,
    citecolor=ipssiVert!80!black,
    bookmarks=true,
    bookmarksopen=true,
    pdfauthor={Mohamed EL AFRIT},
    pdftitle={Mathématiques appliquées à l'IA — IPSSI},
    pdfsubject={Formation Bachelor 2 — IPSSI 2025-2026},
]{hyperref}

% ============================================================================
% IMAGES
% ============================================================================
\usepackage{graphicx}

% ============================================================================
% VARIABLES PARAMÉTRABLES (jour et thème)
% ============================================================================
\newcommand{\leJour}{X}
\newcommand{\leTheme}{Thème}
\newcommand{\jourNumero}[1]{\renewcommand{\leJour}{#1}}
\newcommand{\jourTheme}[1]{\renewcommand{\leTheme}{#1}}

% ============================================================================
% EN-TÊTE ET PIED DE PAGE (fancyhdr)
% ============================================================================
\usepackage{fancyhdr}

\pagestyle{fancy}
\fancyhf{}

% En-tête
\fancyhead[L]{\small\sffamily\textcolor{gray}{IPSSI — Maths appliquées à l'IA}}
\fancyhead[R]{\small\sffamily\textcolor{gray}{Jour \leJour\ — \leTheme}}
\renewcommand{\headrulewidth}{0.4pt}

% Pied de page
\fancyfoot[L]{\footnotesize\textcolor{gray}{Mohamed EL AFRIT\enspace\textbullet\enspace\url{www.mohamedelafrit.com}\enspace\textbullet\enspace CC BY-NC-SA 4.0}}
\fancyfoot[C]{\footnotesize\textcolor{gray}{\thepage}}
\fancyfoot[R]{\footnotesize\textcolor{gray}{2025–2026}}
\renewcommand{\footrulewidth}{0.2pt}

% Style pour la page de garde (pas d'en-tête)
\fancypagestyle{pagegarde}{
    \fancyhf{}
    \renewcommand{\headrulewidth}{0pt}
    \fancyfoot[C]{\footnotesize\textcolor{gray}{\thepage}}
    \renewcommand{\footrulewidth}{0pt}
}

% ============================================================================
% PAGE DE GARDE — \pagegarde{titre}{jour}{theme}
% ============================================================================
\newcommand{\pagegarde}[3]{%
    \thispagestyle{pagegarde}
    \begin{center}
        \vspace*{1.5cm}

        % Logo / Bandeau institution
        {\Large\sffamily\bfseries\textcolor{ipssiBleu}{IPSSI}}\par
        \vspace{2pt}
        {\normalsize\sffamily 2\textsuperscript{e} année Bachelor Informatique}\par

        \vspace{0.8cm}
        {\color{ipssiBleu}\rule{0.6\textwidth}{1.5pt}}\par
        \vspace{0.8cm}

        % Titre du module
        {\LARGE\sffamily\bfseries Mathématiques appliquées\\[4pt] à l'Intelligence Artificielle}\par

        \vspace{1cm}

        % Titre du document
        \begin{tcolorbox}[
            enhanced,
            colback=ipssiBleu!5,
            colframe=ipssiBleu!50,
            boxrule=1pt,
            arc=4pt,
            width=0.85\textwidth,
            halign=center,
            top=10pt, bottom=10pt,
        ]
            {\huge\sffamily\bfseries\textcolor{ipssiBleu!80!black}{#1}}\par
            \vspace{6pt}
            {\Large\sffamily Jour #2 — #3}
        \end{tcolorbox}

        \vspace{1.5cm}

        % Informations auteur
        {\large\sffamily
            \textbf{Auteur :} Mohamed EL AFRIT\\[4pt]
            \url{www.mohamedelafrit.com}\\[8pt]
            \textbf{Année académique :} 2025–2026
        }\par

        \vspace{1.5cm}

        % Licence Creative Commons
        {\color{ipssiBleu}\rule{0.4\textwidth}{0.5pt}}\par
        \vspace{8pt}
        {\small\sffamily
            \textcopyright\ 2025 Mohamed EL AFRIT — IPSSI\\[2pt]
            Ce contenu est distribué sous licence\\[2pt]
            \textbf{Creative Commons BY-NC-SA 4.0}\\[2pt]
            \url{https://creativecommons.org/licenses/by-nc-sa/4.0/deed.fr}
        }\par
    \end{center}
    \newpage
}

% ============================================================================
% NIVEAUX D'EXERCICES
% ============================================================================

% Étoile compatible pdflatex
\newcommand{\starfull}[1]{\textcolor{#1}{$\bigstar$}}
\newcommand{\starempty}{\textcolor{gray!40}{$\bigstar$}}

% Fondamental (1 étoile)
\newcommand{\exofondamental}{%
    \starfull{ipssiVert}\starempty\starempty\enspace%
    {\small\sffamily\textcolor{ipssiVert}{\textbf{Fondamental}}}%
}

% Intermédiaire (2 étoiles)
\newcommand{\exointermediaire}{%
    \starfull{ipssiOrange}\starfull{ipssiOrange}\starempty\enspace%
    {\small\sffamily\textcolor{ipssiOrange}{\textbf{Interm\'{e}diaire}}}%
}

% Avancé (3 étoiles)
\newcommand{\exoavance}{%
    \starfull{ipssiRouge}\starfull{ipssiRouge}\starfull{ipssiRouge}\enspace%
    {\small\sffamily\textcolor{ipssiRouge}{\textbf{Avanc\'{e}}}}%
}

% Commande d'exercice numéroté
\newcounter{exocount}[section]
\newcommand{\exercice}[1]{%
    \stepcounter{exocount}%
    \vspace{8pt}%
    \noindent\textbf{\sffamily Exercice \theexocount}\enspace — \enspace #1\par
    \vspace{4pt}%
}

% ============================================================================
% COMMANDES UTILITAIRES
% ============================================================================

% Séparateur visuel entre sections
\newcommand{\separateur}{%
    \vspace{6pt}
    \begin{center}
        {\color{ipssiBleu!30}\rule{0.3\textwidth}{0.5pt}}
    \end{center}
    \vspace{6pt}
}

% Mot-clé en gras coloré
\newcommand{\motcle}[1]{\textbf{\textcolor{ipssiBleu!80!black}{#1}}}

% Formule encadrée importante
\newcommand{\formule}[1]{%
    \begin{center}
        \fcolorbox{ipssiBleu!50}{ipssiBleu!5}{%
            \quad$\displaystyle #1$\quad%
        }
    \end{center}
}

% ============================================================================
% FIN DU PRÉAMBULE
% ============================================================================
          % Importer ce préambule
%   \jourNumero{1}                  % Définir le numéro du jour (1, 2, 3 ou 4)
%   \jourTheme{Données et représentation mathématique}
%   \begin{document}
%   \pagegarde{Fiche de synthèse}{1}{Données et représentation mathématique}
%   ...
%   \end{document}
% ============================================================================

% --- Classe de document ---
\documentclass[11pt,a4paper]{article}

% ============================================================================
% ENCODAGE ET LANGUE
% ============================================================================
\usepackage[utf8]{inputenc}
\usepackage[T1]{fontenc}
\usepackage[french]{babel}

% ============================================================================
% MATHÉMATIQUES
% ============================================================================
\usepackage{amsmath,amssymb,amsfonts,mathtools}

% Commandes mathématiques personnalisées (conventions du cours)
\newcommand{\vx}{\mathbf{x}}          % vecteur x
\newcommand{\vw}{\mathbf{w}}          % vecteur w (poids)
\newcommand{\vy}{\mathbf{y}}          % vecteur y
\newcommand{\mX}{\mathbf{X}}          % matrice X
\newcommand{\mW}{\mathbf{W}}          % matrice W
\newcommand{\yhat}{\hat{y}}           % prédiction
\newcommand{\pscal}[2]{\langle #1, #2 \rangle} % produit scalaire
\newcommand{\gradJ}{\nabla J(\boldsymbol{\theta})} % gradient de J

% ============================================================================
% MISE EN PAGE
% ============================================================================
\usepackage[margin=2.5cm]{geometry}
\usepackage{setspace}
\setstretch{1.15}                      % interligne légèrement aéré

% ============================================================================
% COULEURS
% ============================================================================
\usepackage[dvipsnames,svgnames,x11names]{xcolor}

% Palette principale du projet
\definecolor{ipssiBleu}{HTML}{3B82F6}       % Définitions
\definecolor{ipssiVert}{HTML}{22C55E}       % À retenir
\definecolor{ipssiOrange}{HTML}{F97316}     % Attention / Piège
\definecolor{ipssiViolet}{HTML}{A855F7}     % Intuition
\definecolor{ipssiCyan}{HTML}{06B6D4}       % Exemples
\definecolor{ipssiGrisClair}{HTML}{F3F4F6}  % Fond code
\definecolor{ipssiGrisFonce}{HTML}{1F2937}  % Texte code
\definecolor{ipssiFondPage}{HTML}{FAFAFA}   % Fond léger page de garde
\definecolor{ipssiRouge}{HTML}{E74C3C}      % Classe 0 (salaire ≤ 45k€)
\definecolor{ipssiVertData}{HTML}{2ECC71}   % Classe 1 (salaire > 45k€)
\definecolor{ipssiBleuDroite}{HTML}{3498DB} % Droite de régression

% ============================================================================
% ENCADRÉS PÉDAGOGIQUES (tcolorbox)
% ============================================================================
\usepackage[most]{tcolorbox}

% Style de base partagé
\tcbset{
    ipssibase/.style={
        enhanced,
        breakable,
        left=4mm,
        right=4mm,
        top=3mm,
        bottom=3mm,
        boxrule=0pt,
        frame hidden,
        borderline west={3pt}{0pt}{#1},
        colback=#1!5,
        colframe=#1,
        fonttitle=\bfseries\sffamily,
        coltitle=#1!80!black,
        attach boxed title to top left={yshift=-2mm, xshift=4mm},
        boxed title style={
            colback=#1!15,
            colframe=#1,
            boxrule=0.5pt,
            arc=2pt,
        },
        arc=2pt,
        outer arc=2pt,
    }
}

% Icônes TikZ pour les encadrés (compatible pdflatex, pas d'emoji Unicode)
\newcommand{\icoDefinition}{%
    \tikz[baseline=-0.3em]{\draw[ipssiBleu, thick] (0,0) -- (0.3,0) -- (0.3,0.35) -- (0,0.35) -- cycle;
    \draw[ipssiBleu, thick] (0.05,0.1) -- (0.25,0.1); \draw[ipssiBleu, thick] (0.05,0.2) -- (0.2,0.2);}%
}
\newcommand{\icoRetenir}{%
    \tikz[baseline=-0.2em]{\draw[ipssiVert, thick, fill=ipssiVert!20] (0.15,0) -- (0.3,0.3) -- (0,0.3) -- cycle;}%
}
\newcommand{\icoAttention}{%
    \tikz[baseline=-0.2em]{\draw[ipssiOrange, thick] (0.15,0.35) -- (0,0) -- (0.3,0) -- cycle;
    \node[ipssiOrange, font=\tiny\bfseries] at (0.15,0.12) {!};}%
}
\newcommand{\icoIntuition}{%
    \tikz[baseline=-0.2em]{\draw[ipssiViolet, thick, fill=ipssiViolet!15]
    (0.15,0.35) -- (0.08,0.2) -- (0.08,0.15) -- (0.22,0.15) -- (0.22,0.2) -- cycle;
    \draw[ipssiViolet, thick] (0.11,0.12) -- (0.11,0.05); \draw[ipssiViolet, thick] (0.19,0.12) -- (0.19,0.05);}%
}
\newcommand{\icoExemple}{%
    \tikz[baseline=-0.2em]{\draw[ipssiCyan, thick] (0,0) rectangle (0.3,0.3);
    \draw[ipssiCyan, thick] (0,0.15) -- (0.3,0.15); \draw[ipssiCyan, thick] (0.15,0) -- (0.15,0.3);}%
}
\newcommand{\icoPython}{%
    \tikz[baseline=-0.2em]{\node[font=\small\bfseries\ttfamily, ipssiVert!70!black] at (0,0) {\textgreater\textgreater};}%
}

% Définition (bleu)
\newtcolorbox{definitionbox}[1][D\'{e}finition]{
    ipssibase=ipssiBleu,
    title={\icoDefinition\; #1},
}

% À retenir (vert)
\newtcolorbox{retenirbox}[1][\`{A} retenir]{
    ipssibase=ipssiVert,
    title={\icoRetenir\; #1},
}

% Attention / Piège (orange)
\newtcolorbox{attentionbox}[1][Attention]{
    ipssibase=ipssiOrange,
    title={\icoAttention\; #1},
}

% Intuition (violet)
\newtcolorbox{intuitionbox}[1][Intuition]{
    ipssibase=ipssiViolet,
    title={\icoIntuition\; #1},
}

% Exemple (cyan)
\newtcolorbox{exemplebox}[1][Exemple]{
    ipssibase=ipssiCyan,
    title={\icoExemple\; #1},
}

% Code Python (style terminal)
\newtcolorbox{pythonbox}[1][Code Python]{
    enhanced,
    breakable,
    left=4mm, right=4mm, top=3mm, bottom=3mm,
    boxrule=0.5pt,
    colback=ipssiGrisClair,
    colframe=ipssiGrisFonce!40,
    fonttitle=\bfseries\ttfamily\small,
    coltitle=ipssiGrisFonce,
    title={\icoPython\; #1},
    attach boxed title to top left={yshift=-2mm, xshift=4mm},
    boxed title style={colback=ipssiGrisClair, colframe=ipssiGrisFonce!40, boxrule=0.5pt, arc=2pt},
    arc=2pt, outer arc=2pt,
}

% ============================================================================
% CODE PYTHON (listings)
% ============================================================================
\usepackage{listings}

\lstdefinestyle{python}{
    language=Python,
    basicstyle=\ttfamily\small,
    keywordstyle=\color{ipssiBleu}\bfseries,
    stringstyle=\color{ipssiVert!70!black},
    commentstyle=\color{gray}\itshape,
    numberstyle=\tiny\color{gray},
    numbers=left,
    numbersep=8pt,
    backgroundcolor=\color{ipssiGrisClair},
    frame=none,
    rulecolor=\color{ipssiGrisFonce!20},
    breaklines=true,
    breakatwhitespace=true,
    tabsize=4,
    showstringspaces=false,
    captionpos=b,
    aboveskip=8pt,
    belowskip=8pt,
    xleftmargin=10pt,
    framexleftmargin=10pt,
    morekeywords={as, with, True, False, None, self, np, plt, import, from},
    literate=
        {é}{{\'e}}1 {è}{{\`e}}1 {ê}{{\^e}}1 {ë}{{\"e}}1
        {à}{{\`a}}1 {â}{{\^a}}1
        {ù}{{\`u}}1 {û}{{\^u}}1
        {ô}{{\^o}}1
        {î}{{\^i}}1 {ï}{{\"i}}1
        {ç}{{\c{c}}}1
        {É}{{\'E}}1 {È}{{\`E}}1
        {→}{{$\rightarrow$}}1
        {≈}{{$\approx$}}1
        {≤}{{$\leq$}}1
        {≥}{{$\geq$}}1
        {★}{{$\bigstar$}}1
        {ŷ}{{$\hat{y}$}}1
        {·}{{$\cdot$}}1,
}

\lstset{style=python}

% Style pour le code inline
\newcommand{\pyinline}[1]{\lstinline[style=python,basicstyle=\ttfamily\small]{#1}}

% ============================================================================
% GRAPHIQUES
% ============================================================================
\usepackage{tikz}
\usepackage{pgfplots}
\pgfplotsset{compat=1.18}
\usetikzlibrary{arrows.meta, calc, positioning, shapes.geometric, decorations.pathreplacing}

% ============================================================================
% TABLEAUX
% ============================================================================
\usepackage{booktabs}
\usepackage{tabularx}
\usepackage{multirow}

% ============================================================================
% LISTES
% ============================================================================
\usepackage{enumitem}
\setlist[itemize]{itemsep=2pt, parsep=0pt, topsep=4pt}
\setlist[enumerate]{itemsep=2pt, parsep=0pt, topsep=4pt}

% ============================================================================
% HYPERLIENS
% ============================================================================
\usepackage[
    colorlinks=true,
    linkcolor=ipssiBleu!80!black,
    urlcolor=ipssiBleu!80!black,
    citecolor=ipssiVert!80!black,
    bookmarks=true,
    bookmarksopen=true,
    pdfauthor={Mohamed EL AFRIT},
    pdftitle={Mathématiques appliquées à l'IA — IPSSI},
    pdfsubject={Formation Bachelor 2 — IPSSI 2025-2026},
]{hyperref}

% ============================================================================
% IMAGES
% ============================================================================
\usepackage{graphicx}

% ============================================================================
% VARIABLES PARAMÉTRABLES (jour et thème)
% ============================================================================
\newcommand{\leJour}{X}
\newcommand{\leTheme}{Thème}
\newcommand{\jourNumero}[1]{\renewcommand{\leJour}{#1}}
\newcommand{\jourTheme}[1]{\renewcommand{\leTheme}{#1}}

% ============================================================================
% EN-TÊTE ET PIED DE PAGE (fancyhdr)
% ============================================================================
\usepackage{fancyhdr}

\pagestyle{fancy}
\fancyhf{}

% En-tête
\fancyhead[L]{\small\sffamily\textcolor{gray}{IPSSI — Maths appliquées à l'IA}}
\fancyhead[R]{\small\sffamily\textcolor{gray}{Jour \leJour\ — \leTheme}}
\renewcommand{\headrulewidth}{0.4pt}

% Pied de page
\fancyfoot[L]{\footnotesize\textcolor{gray}{Mohamed EL AFRIT\enspace\textbullet\enspace\url{www.mohamedelafrit.com}\enspace\textbullet\enspace CC BY-NC-SA 4.0}}
\fancyfoot[C]{\footnotesize\textcolor{gray}{\thepage}}
\fancyfoot[R]{\footnotesize\textcolor{gray}{2025–2026}}
\renewcommand{\footrulewidth}{0.2pt}

% Style pour la page de garde (pas d'en-tête)
\fancypagestyle{pagegarde}{
    \fancyhf{}
    \renewcommand{\headrulewidth}{0pt}
    \fancyfoot[C]{\footnotesize\textcolor{gray}{\thepage}}
    \renewcommand{\footrulewidth}{0pt}
}

% ============================================================================
% PAGE DE GARDE — \pagegarde{titre}{jour}{theme}
% ============================================================================
\newcommand{\pagegarde}[3]{%
    \thispagestyle{pagegarde}
    \begin{center}
        \vspace*{1.5cm}

        % Logo / Bandeau institution
        {\Large\sffamily\bfseries\textcolor{ipssiBleu}{IPSSI}}\par
        \vspace{2pt}
        {\normalsize\sffamily 2\textsuperscript{e} année Bachelor Informatique}\par

        \vspace{0.8cm}
        {\color{ipssiBleu}\rule{0.6\textwidth}{1.5pt}}\par
        \vspace{0.8cm}

        % Titre du module
        {\LARGE\sffamily\bfseries Mathématiques appliquées\\[4pt] à l'Intelligence Artificielle}\par

        \vspace{1cm}

        % Titre du document
        \begin{tcolorbox}[
            enhanced,
            colback=ipssiBleu!5,
            colframe=ipssiBleu!50,
            boxrule=1pt,
            arc=4pt,
            width=0.85\textwidth,
            halign=center,
            top=10pt, bottom=10pt,
        ]
            {\huge\sffamily\bfseries\textcolor{ipssiBleu!80!black}{#1}}\par
            \vspace{6pt}
            {\Large\sffamily Jour #2 — #3}
        \end{tcolorbox}

        \vspace{1.5cm}

        % Informations auteur
        {\large\sffamily
            \textbf{Auteur :} Mohamed EL AFRIT\\[4pt]
            \url{www.mohamedelafrit.com}\\[8pt]
            \textbf{Année académique :} 2025–2026
        }\par

        \vspace{1.5cm}

        % Licence Creative Commons
        {\color{ipssiBleu}\rule{0.4\textwidth}{0.5pt}}\par
        \vspace{8pt}
        {\small\sffamily
            \textcopyright\ 2025 Mohamed EL AFRIT — IPSSI\\[2pt]
            Ce contenu est distribué sous licence\\[2pt]
            \textbf{Creative Commons BY-NC-SA 4.0}\\[2pt]
            \url{https://creativecommons.org/licenses/by-nc-sa/4.0/deed.fr}
        }\par
    \end{center}
    \newpage
}

% ============================================================================
% NIVEAUX D'EXERCICES
% ============================================================================

% Étoile compatible pdflatex
\newcommand{\starfull}[1]{\textcolor{#1}{$\bigstar$}}
\newcommand{\starempty}{\textcolor{gray!40}{$\bigstar$}}

% Fondamental (1 étoile)
\newcommand{\exofondamental}{%
    \starfull{ipssiVert}\starempty\starempty\enspace%
    {\small\sffamily\textcolor{ipssiVert}{\textbf{Fondamental}}}%
}

% Intermédiaire (2 étoiles)
\newcommand{\exointermediaire}{%
    \starfull{ipssiOrange}\starfull{ipssiOrange}\starempty\enspace%
    {\small\sffamily\textcolor{ipssiOrange}{\textbf{Interm\'{e}diaire}}}%
}

% Avancé (3 étoiles)
\newcommand{\exoavance}{%
    \starfull{ipssiRouge}\starfull{ipssiRouge}\starfull{ipssiRouge}\enspace%
    {\small\sffamily\textcolor{ipssiRouge}{\textbf{Avanc\'{e}}}}%
}

% Commande d'exercice numéroté
\newcounter{exocount}[section]
\newcommand{\exercice}[1]{%
    \stepcounter{exocount}%
    \vspace{8pt}%
    \noindent\textbf{\sffamily Exercice \theexocount}\enspace — \enspace #1\par
    \vspace{4pt}%
}

% ============================================================================
% COMMANDES UTILITAIRES
% ============================================================================

% Séparateur visuel entre sections
\newcommand{\separateur}{%
    \vspace{6pt}
    \begin{center}
        {\color{ipssiBleu!30}\rule{0.3\textwidth}{0.5pt}}
    \end{center}
    \vspace{6pt}
}

% Mot-clé en gras coloré
\newcommand{\motcle}[1]{\textbf{\textcolor{ipssiBleu!80!black}{#1}}}

% Formule encadrée importante
\newcommand{\formule}[1]{%
    \begin{center}
        \fcolorbox{ipssiBleu!50}{ipssiBleu!5}{%
            \quad$\displaystyle #1$\quad%
        }
    \end{center}
}

% ============================================================================
% FIN DU PRÉAMBULE
% ============================================================================
          % Importer ce préambule
%   \jourNumero{1}                  % Définir le numéro du jour (1, 2, 3 ou 4)
%   \jourTheme{Données et représentation mathématique}
%   \begin{document}
%   \pagegarde{Fiche de synthèse}{1}{Données et représentation mathématique}
%   ...
%   \end{document}
% ============================================================================

% --- Classe de document ---
\documentclass[11pt,a4paper]{article}

% ============================================================================
% ENCODAGE ET LANGUE
% ============================================================================
\usepackage[utf8]{inputenc}
\usepackage[T1]{fontenc}
\usepackage[french]{babel}

% ============================================================================
% MATHÉMATIQUES
% ============================================================================
\usepackage{amsmath,amssymb,amsfonts,mathtools}

% Commandes mathématiques personnalisées (conventions du cours)
\newcommand{\vx}{\mathbf{x}}          % vecteur x
\newcommand{\vw}{\mathbf{w}}          % vecteur w (poids)
\newcommand{\vy}{\mathbf{y}}          % vecteur y
\newcommand{\mX}{\mathbf{X}}          % matrice X
\newcommand{\mW}{\mathbf{W}}          % matrice W
\newcommand{\yhat}{\hat{y}}           % prédiction
\newcommand{\pscal}[2]{\langle #1, #2 \rangle} % produit scalaire
\newcommand{\gradJ}{\nabla J(\boldsymbol{\theta})} % gradient de J

% ============================================================================
% MISE EN PAGE
% ============================================================================
\usepackage[margin=2.5cm]{geometry}
\usepackage{setspace}
\setstretch{1.15}                      % interligne légèrement aéré

% ============================================================================
% COULEURS
% ============================================================================
\usepackage[dvipsnames,svgnames,x11names]{xcolor}

% Palette principale du projet
\definecolor{ipssiBleu}{HTML}{3B82F6}       % Définitions
\definecolor{ipssiVert}{HTML}{22C55E}       % À retenir
\definecolor{ipssiOrange}{HTML}{F97316}     % Attention / Piège
\definecolor{ipssiViolet}{HTML}{A855F7}     % Intuition
\definecolor{ipssiCyan}{HTML}{06B6D4}       % Exemples
\definecolor{ipssiGrisClair}{HTML}{F3F4F6}  % Fond code
\definecolor{ipssiGrisFonce}{HTML}{1F2937}  % Texte code
\definecolor{ipssiFondPage}{HTML}{FAFAFA}   % Fond léger page de garde
\definecolor{ipssiRouge}{HTML}{E74C3C}      % Classe 0 (salaire ≤ 45k€)
\definecolor{ipssiVertData}{HTML}{2ECC71}   % Classe 1 (salaire > 45k€)
\definecolor{ipssiBleuDroite}{HTML}{3498DB} % Droite de régression

% ============================================================================
% ENCADRÉS PÉDAGOGIQUES (tcolorbox)
% ============================================================================
\usepackage[most]{tcolorbox}

% Style de base partagé
\tcbset{
    ipssibase/.style={
        enhanced,
        breakable,
        left=4mm,
        right=4mm,
        top=3mm,
        bottom=3mm,
        boxrule=0pt,
        frame hidden,
        borderline west={3pt}{0pt}{#1},
        colback=#1!5,
        colframe=#1,
        fonttitle=\bfseries\sffamily,
        coltitle=#1!80!black,
        attach boxed title to top left={yshift=-2mm, xshift=4mm},
        boxed title style={
            colback=#1!15,
            colframe=#1,
            boxrule=0.5pt,
            arc=2pt,
        },
        arc=2pt,
        outer arc=2pt,
    }
}

% Icônes TikZ pour les encadrés (compatible pdflatex, pas d'emoji Unicode)
\newcommand{\icoDefinition}{%
    \tikz[baseline=-0.3em]{\draw[ipssiBleu, thick] (0,0) -- (0.3,0) -- (0.3,0.35) -- (0,0.35) -- cycle;
    \draw[ipssiBleu, thick] (0.05,0.1) -- (0.25,0.1); \draw[ipssiBleu, thick] (0.05,0.2) -- (0.2,0.2);}%
}
\newcommand{\icoRetenir}{%
    \tikz[baseline=-0.2em]{\draw[ipssiVert, thick, fill=ipssiVert!20] (0.15,0) -- (0.3,0.3) -- (0,0.3) -- cycle;}%
}
\newcommand{\icoAttention}{%
    \tikz[baseline=-0.2em]{\draw[ipssiOrange, thick] (0.15,0.35) -- (0,0) -- (0.3,0) -- cycle;
    \node[ipssiOrange, font=\tiny\bfseries] at (0.15,0.12) {!};}%
}
\newcommand{\icoIntuition}{%
    \tikz[baseline=-0.2em]{\draw[ipssiViolet, thick, fill=ipssiViolet!15]
    (0.15,0.35) -- (0.08,0.2) -- (0.08,0.15) -- (0.22,0.15) -- (0.22,0.2) -- cycle;
    \draw[ipssiViolet, thick] (0.11,0.12) -- (0.11,0.05); \draw[ipssiViolet, thick] (0.19,0.12) -- (0.19,0.05);}%
}
\newcommand{\icoExemple}{%
    \tikz[baseline=-0.2em]{\draw[ipssiCyan, thick] (0,0) rectangle (0.3,0.3);
    \draw[ipssiCyan, thick] (0,0.15) -- (0.3,0.15); \draw[ipssiCyan, thick] (0.15,0) -- (0.15,0.3);}%
}
\newcommand{\icoPython}{%
    \tikz[baseline=-0.2em]{\node[font=\small\bfseries\ttfamily, ipssiVert!70!black] at (0,0) {\textgreater\textgreater};}%
}

% Définition (bleu)
\newtcolorbox{definitionbox}[1][D\'{e}finition]{
    ipssibase=ipssiBleu,
    title={\icoDefinition\; #1},
}

% À retenir (vert)
\newtcolorbox{retenirbox}[1][\`{A} retenir]{
    ipssibase=ipssiVert,
    title={\icoRetenir\; #1},
}

% Attention / Piège (orange)
\newtcolorbox{attentionbox}[1][Attention]{
    ipssibase=ipssiOrange,
    title={\icoAttention\; #1},
}

% Intuition (violet)
\newtcolorbox{intuitionbox}[1][Intuition]{
    ipssibase=ipssiViolet,
    title={\icoIntuition\; #1},
}

% Exemple (cyan)
\newtcolorbox{exemplebox}[1][Exemple]{
    ipssibase=ipssiCyan,
    title={\icoExemple\; #1},
}

% Code Python (style terminal)
\newtcolorbox{pythonbox}[1][Code Python]{
    enhanced,
    breakable,
    left=4mm, right=4mm, top=3mm, bottom=3mm,
    boxrule=0.5pt,
    colback=ipssiGrisClair,
    colframe=ipssiGrisFonce!40,
    fonttitle=\bfseries\ttfamily\small,
    coltitle=ipssiGrisFonce,
    title={\icoPython\; #1},
    attach boxed title to top left={yshift=-2mm, xshift=4mm},
    boxed title style={colback=ipssiGrisClair, colframe=ipssiGrisFonce!40, boxrule=0.5pt, arc=2pt},
    arc=2pt, outer arc=2pt,
}

% ============================================================================
% CODE PYTHON (listings)
% ============================================================================
\usepackage{listings}

\lstdefinestyle{python}{
    language=Python,
    basicstyle=\ttfamily\small,
    keywordstyle=\color{ipssiBleu}\bfseries,
    stringstyle=\color{ipssiVert!70!black},
    commentstyle=\color{gray}\itshape,
    numberstyle=\tiny\color{gray},
    numbers=left,
    numbersep=8pt,
    backgroundcolor=\color{ipssiGrisClair},
    frame=none,
    rulecolor=\color{ipssiGrisFonce!20},
    breaklines=true,
    breakatwhitespace=true,
    tabsize=4,
    showstringspaces=false,
    captionpos=b,
    aboveskip=8pt,
    belowskip=8pt,
    xleftmargin=10pt,
    framexleftmargin=10pt,
    morekeywords={as, with, True, False, None, self, np, plt, import, from},
    literate=
        {é}{{\'e}}1 {è}{{\`e}}1 {ê}{{\^e}}1 {ë}{{\"e}}1
        {à}{{\`a}}1 {â}{{\^a}}1
        {ù}{{\`u}}1 {û}{{\^u}}1
        {ô}{{\^o}}1
        {î}{{\^i}}1 {ï}{{\"i}}1
        {ç}{{\c{c}}}1
        {É}{{\'E}}1 {È}{{\`E}}1
        {→}{{$\rightarrow$}}1
        {≈}{{$\approx$}}1
        {≤}{{$\leq$}}1
        {≥}{{$\geq$}}1
        {★}{{$\bigstar$}}1
        {ŷ}{{$\hat{y}$}}1
        {·}{{$\cdot$}}1,
}

\lstset{style=python}

% Style pour le code inline
\newcommand{\pyinline}[1]{\lstinline[style=python,basicstyle=\ttfamily\small]{#1}}

% ============================================================================
% GRAPHIQUES
% ============================================================================
\usepackage{tikz}
\usepackage{pgfplots}
\pgfplotsset{compat=1.18}
\usetikzlibrary{arrows.meta, calc, positioning, shapes.geometric, decorations.pathreplacing}

% ============================================================================
% TABLEAUX
% ============================================================================
\usepackage{booktabs}
\usepackage{tabularx}
\usepackage{multirow}

% ============================================================================
% LISTES
% ============================================================================
\usepackage{enumitem}
\setlist[itemize]{itemsep=2pt, parsep=0pt, topsep=4pt}
\setlist[enumerate]{itemsep=2pt, parsep=0pt, topsep=4pt}

% ============================================================================
% HYPERLIENS
% ============================================================================
\usepackage[
    colorlinks=true,
    linkcolor=ipssiBleu!80!black,
    urlcolor=ipssiBleu!80!black,
    citecolor=ipssiVert!80!black,
    bookmarks=true,
    bookmarksopen=true,
    pdfauthor={Mohamed EL AFRIT},
    pdftitle={Mathématiques appliquées à l'IA — IPSSI},
    pdfsubject={Formation Bachelor 2 — IPSSI 2025-2026},
]{hyperref}

% ============================================================================
% IMAGES
% ============================================================================
\usepackage{graphicx}

% ============================================================================
% VARIABLES PARAMÉTRABLES (jour et thème)
% ============================================================================
\newcommand{\leJour}{X}
\newcommand{\leTheme}{Thème}
\newcommand{\jourNumero}[1]{\renewcommand{\leJour}{#1}}
\newcommand{\jourTheme}[1]{\renewcommand{\leTheme}{#1}}

% ============================================================================
% EN-TÊTE ET PIED DE PAGE (fancyhdr)
% ============================================================================
\usepackage{fancyhdr}

\pagestyle{fancy}
\fancyhf{}

% En-tête
\fancyhead[L]{\small\sffamily\textcolor{gray}{IPSSI — Maths appliquées à l'IA}}
\fancyhead[R]{\small\sffamily\textcolor{gray}{Jour \leJour\ — \leTheme}}
\renewcommand{\headrulewidth}{0.4pt}

% Pied de page
\fancyfoot[L]{\footnotesize\textcolor{gray}{Mohamed EL AFRIT\enspace\textbullet\enspace\url{www.mohamedelafrit.com}\enspace\textbullet\enspace CC BY-NC-SA 4.0}}
\fancyfoot[C]{\footnotesize\textcolor{gray}{\thepage}}
\fancyfoot[R]{\footnotesize\textcolor{gray}{2025–2026}}
\renewcommand{\footrulewidth}{0.2pt}

% Style pour la page de garde (pas d'en-tête)
\fancypagestyle{pagegarde}{
    \fancyhf{}
    \renewcommand{\headrulewidth}{0pt}
    \fancyfoot[C]{\footnotesize\textcolor{gray}{\thepage}}
    \renewcommand{\footrulewidth}{0pt}
}

% ============================================================================
% PAGE DE GARDE — \pagegarde{titre}{jour}{theme}
% ============================================================================
\newcommand{\pagegarde}[3]{%
    \thispagestyle{pagegarde}
    \begin{center}
        \vspace*{1.5cm}

        % Logo / Bandeau institution
        {\Large\sffamily\bfseries\textcolor{ipssiBleu}{IPSSI}}\par
        \vspace{2pt}
        {\normalsize\sffamily 2\textsuperscript{e} année Bachelor Informatique}\par

        \vspace{0.8cm}
        {\color{ipssiBleu}\rule{0.6\textwidth}{1.5pt}}\par
        \vspace{0.8cm}

        % Titre du module
        {\LARGE\sffamily\bfseries Mathématiques appliquées\\[4pt] à l'Intelligence Artificielle}\par

        \vspace{1cm}

        % Titre du document
        \begin{tcolorbox}[
            enhanced,
            colback=ipssiBleu!5,
            colframe=ipssiBleu!50,
            boxrule=1pt,
            arc=4pt,
            width=0.85\textwidth,
            halign=center,
            top=10pt, bottom=10pt,
        ]
            {\huge\sffamily\bfseries\textcolor{ipssiBleu!80!black}{#1}}\par
            \vspace{6pt}
            {\Large\sffamily Jour #2 — #3}
        \end{tcolorbox}

        \vspace{1.5cm}

        % Informations auteur
        {\large\sffamily
            \textbf{Auteur :} Mohamed EL AFRIT\\[4pt]
            \url{www.mohamedelafrit.com}\\[8pt]
            \textbf{Année académique :} 2025–2026
        }\par

        \vspace{1.5cm}

        % Licence Creative Commons
        {\color{ipssiBleu}\rule{0.4\textwidth}{0.5pt}}\par
        \vspace{8pt}
        {\small\sffamily
            \textcopyright\ 2025 Mohamed EL AFRIT — IPSSI\\[2pt]
            Ce contenu est distribué sous licence\\[2pt]
            \textbf{Creative Commons BY-NC-SA 4.0}\\[2pt]
            \url{https://creativecommons.org/licenses/by-nc-sa/4.0/deed.fr}
        }\par
    \end{center}
    \newpage
}

% ============================================================================
% NIVEAUX D'EXERCICES
% ============================================================================

% Étoile compatible pdflatex
\newcommand{\starfull}[1]{\textcolor{#1}{$\bigstar$}}
\newcommand{\starempty}{\textcolor{gray!40}{$\bigstar$}}

% Fondamental (1 étoile)
\newcommand{\exofondamental}{%
    \starfull{ipssiVert}\starempty\starempty\enspace%
    {\small\sffamily\textcolor{ipssiVert}{\textbf{Fondamental}}}%
}

% Intermédiaire (2 étoiles)
\newcommand{\exointermediaire}{%
    \starfull{ipssiOrange}\starfull{ipssiOrange}\starempty\enspace%
    {\small\sffamily\textcolor{ipssiOrange}{\textbf{Interm\'{e}diaire}}}%
}

% Avancé (3 étoiles)
\newcommand{\exoavance}{%
    \starfull{ipssiRouge}\starfull{ipssiRouge}\starfull{ipssiRouge}\enspace%
    {\small\sffamily\textcolor{ipssiRouge}{\textbf{Avanc\'{e}}}}%
}

% Commande d'exercice numéroté
\newcounter{exocount}[section]
\newcommand{\exercice}[1]{%
    \stepcounter{exocount}%
    \vspace{8pt}%
    \noindent\textbf{\sffamily Exercice \theexocount}\enspace — \enspace #1\par
    \vspace{4pt}%
}

% ============================================================================
% COMMANDES UTILITAIRES
% ============================================================================

% Séparateur visuel entre sections
\newcommand{\separateur}{%
    \vspace{6pt}
    \begin{center}
        {\color{ipssiBleu!30}\rule{0.3\textwidth}{0.5pt}}
    \end{center}
    \vspace{6pt}
}

% Mot-clé en gras coloré
\newcommand{\motcle}[1]{\textbf{\textcolor{ipssiBleu!80!black}{#1}}}

% Formule encadrée importante
\newcommand{\formule}[1]{%
    \begin{center}
        \fcolorbox{ipssiBleu!50}{ipssiBleu!5}{%
            \quad$\displaystyle #1$\quad%
        }
    \end{center}
}

% ============================================================================
% FIN DU PRÉAMBULE
% ============================================================================


\jourNumero{1}
\jourTheme{Repr\'{e}sentation}

\begin{document}

\pagegarde{Fiche d'exercices}{1}{Donn\'{e}es et repr\'{e}sentation math\'{e}matique}

% ── Compacter pour la densité ──────────────────────────────────────────────
\setstretch{1.08}
\setlength{\parskip}{3pt}
\setlength{\abovedisplayskip}{6pt}
\setlength{\belowdisplayskip}{6pt}

% ══════════════════════════════════════════════════════════════════════════
% RAPPEL DU DATASET FIL ROUGE
% ══════════════════════════════════════════════════════════════════════════
\begin{tcolorbox}[
    enhanced, colback=ipssiCyan!5, colframe=ipssiCyan!50,
    boxrule=0.8pt, arc=4pt, title={\sffamily\bfseries Rappel --- Dataset fil rouge},
    coltitle=white, colbacktitle=ipssiCyan!70,
]
\small
Tous les exercices utilisent le dataset \textbf{Salaire d\'{e}veloppeur} (30 observations, 5 features).
Colonnes : \texttt{experience} (0--20), \texttt{nb\_langages} (1--8), \texttt{niveau\_etudes} (1--5),
\texttt{taille\_entreprise} (15--5000), \texttt{remote} (0--100), \texttt{salaire} (25--72 k\texteuro).

\medskip
\begin{center}
\renewcommand{\arraystretch}{1.1}
\begin{tabular}{c|ccccc|c}
    \toprule
    \textbf{\#} & \textbf{Exp} & \textbf{Lang} & \textbf{\'{E}tudes} & \textbf{Entrep.} & \textbf{Remote} & \textbf{Salaire} \\
    \midrule
    0  & 0  & 2 & 3 & 150  & 80  & 28.0 \\
    1  & 1  & 1 & 2 & 50   & 100 & 26.5 \\
    2  & 1  & 3 & 4 & 800  & 40  & 35.0 \\
    3  & 2  & 2 & 3 & 200  & 60  & 32.0 \\
    4  & 2  & 4 & 4 & 3000 & 20  & 38.0 \\
    5  & 0  & 1 & 1 & 30   & 50  & 25.0 \\
    \bottomrule
\end{tabular}
\end{center}

R\'{e}gression simple (exp\'{e}rience $\to$ salaire) :
$\yhat = 2{,}07 \times \texttt{experience} + 31{,}27$ \quad (MSE $= 20{,}46$).
\end{tcolorbox}

% ══════════════════════════════════════════════════════════════════════════
\section*{Exercices crayon-papier --- Niveau fondamental}
% ══════════════════════════════════════════════════════════════════════════

% ──────────────────────────────────────────────────────────────────────────
\exercice{\exofondamental\quad --- Vecteurs et dimensions}

\begin{enumerate}[label=\alph*)]
    \item \'{E}crivez le vecteur de features $\vx_0$ du d\'{e}veloppeur \#0 du dataset.
        Quelle est sa dimension~$p$ ?

    \item \'{E}crivez le vecteur $\vx_4$ du d\'{e}veloppeur \#4. \'{E}crivez sa transpos\'{e}e $\vx_4^T$.

    \item Un nouveau d\'{e}veloppeur a le profil suivant : 6 ans d'exp\'{e}rience,
        4 langages, Bac+5, entreprise de 1200 employ\'{e}s, 50\% de remote.
        \'{E}crivez son vecteur $\vx$ et pr\'{e}cisez $\vx \in \mathbb{R}^{?}$.

    \item Peut-on calculer $\vx_0 + \vx_4$ ? Et $\vx_0 + \begin{pmatrix} 1 \\ 2 \\ 3 \end{pmatrix}$ ?
        Justifiez.
\end{enumerate}

% ──────────────────────────────────────────────────────────────────────────
\exercice{\exofondamental\quad --- Produit scalaire \`{a} la main}

\begin{enumerate}[label=\alph*)]
    \item Calculez le produit scalaire $\pscal{\vx}{\vw}$ avec :
    \[
        \vx = \begin{pmatrix} 3 \\ 5 \\ 2 \end{pmatrix}
        \qquad
        \vw = \begin{pmatrix} 0.5 \\ 0.3 \\ 0.8 \end{pmatrix}
    \]
    D\'{e}taillez le calcul terme \`{a} terme.

    \item Calculez $\pscal{\vx_0}{\vw}$ en utilisant le d\'{e}veloppeur \#0 du dataset
        et le vecteur de poids $\vw = \begin{pmatrix} 2.0 \\ 1.5 \\ 3.0 \\ 0.005 \\ 0.02 \end{pmatrix}$.
        Ajoutez un biais $b = 15$ pour obtenir la pr\'{e}diction $\yhat_0$.
        Comparez avec le vrai salaire $y_0 = 28{,}0$ k\texteuro.

    \item \textbf{Pi\`{e}ge :} peut-on calculer le produit scalaire suivant ?
    \[
        \left\langle
        \begin{pmatrix} 1 \\ 2 \\ 3 \\ 4 \\ 5 \end{pmatrix},\;
        \begin{pmatrix} 0.1 \\ 0.2 \\ 0.3 \end{pmatrix}
        \right\rangle
    \]
    Expliquez pourquoi.
\end{enumerate}

% ──────────────────────────────────────────────────────────────────────────
\exercice{\exofondamental\quad --- Multiplication matrice-vecteur}

Soit la matrice $\mX$ et le vecteur $\vw$ suivants :
\[
    \mX = \begin{pmatrix}
        1 & 3 \\
        4 & 2 \\
        0 & 5
    \end{pmatrix}
    \qquad
    \vw = \begin{pmatrix} 2 \\ 0.5 \end{pmatrix}
\]

\begin{enumerate}[label=\alph*)]
    \item Quelle est la dimension de $\mX$ ? De $\vw$ ?
        Quelle sera la dimension du r\'{e}sultat $\hat{\vy} = \mX\vw$ ?

    \item Calculez $\hat{\vy} = \mX\vw$ en d\'{e}taillant chaque produit scalaire
        (ligne par ligne).

    \item V\'{e}rifiez que $\yhat_1 = \vx_1^T \vw$ o\`{u} $\vx_1^T$ est la premi\`{e}re
        ligne de $\mX$.
\end{enumerate}

% ──────────────────────────────────────────────────────────────────────────
\exercice{\exofondamental\quad --- Calcul du MSE}

Un mod\`{e}le a produit les pr\'{e}dictions suivantes pour 5 d\'{e}veloppeurs :

\begin{center}
\renewcommand{\arraystretch}{1.2}
\begin{tabular}{c|c|c}
    \toprule
    D\'{e}veloppeur $i$ & Salaire r\'{e}el $y_i$ (k\texteuro) & Pr\'{e}diction $\yhat_i$ (k\texteuro) \\
    \midrule
    1 & 28.0 & 31.3 \\
    2 & 35.0 & 33.3 \\
    3 & 38.0 & 35.4 \\
    4 & 44.0 & 41.6 \\
    5 & 52.0 & 49.9 \\
    \bottomrule
\end{tabular}
\end{center}

\begin{enumerate}[label=\alph*)]
    \item Calculez l'erreur (r\'{e}sidu) $e_i = y_i - \yhat_i$ pour chaque observation.

    \item Calculez l'erreur au carr\'{e} $e_i^2$ pour chaque observation.

    \item Calculez le MSE : $\text{MSE} = \frac{1}{n}\sum_{i=1}^{n} e_i^2$.

    \item L'erreur typique du mod\`{e}le est $\sqrt{\text{MSE}}$.
        Calculez-la et interpr\'{e}tez : est-ce que le mod\`{e}le pr\'{e}dit bien ?
\end{enumerate}

% ══════════════════════════════════════════════════════════════════════════
\section*{Exercices crayon-papier --- Niveau interm\'{e}diaire}
% ══════════════════════════════════════════════════════════════════════════

% ──────────────────────────────────────────────────────────────────────────
\exercice{\exointermediaire\quad --- Interpr\'{e}ter les poids d'un mod\`{e}le}

On consid\`{e}re le mod\`{e}le lin\'{e}aire suivant pour pr\'{e}dire le salaire (en k\texteuro) :
\[
    \yhat = 25 + 3 \times \texttt{experience} + 1{,}5 \times \texttt{nb\_langages}
\]

\begin{enumerate}[label=\alph*)]
    \item Identifiez $w_0$ (biais), $w_1$ et $w_2$.
        Que repr\'{e}sente concr\`{e}tement chacun de ces param\`{e}tres ?

    \item Quel salaire ce mod\`{e}le pr\'{e}dit-il pour un d\'{e}veloppeur avec 0 ann\'{e}e
        d'exp\'{e}rience et 1 langage ? Est-ce r\'{e}aliste ?

    \item Quel est l'impact d'une ann\'{e}e d'exp\'{e}rience suppl\'{e}mentaire sur le salaire
        pr\'{e}dit ? D'un langage suppl\'{e}mentaire ?

    \item Si un d\'{e}veloppeur a 5 ans d'exp\'{e}rience et ma\^{i}trise 4 langages,
        quel salaire le mod\`{e}le pr\'{e}dit-il ?
        Comparez avec le d\'{e}veloppeur \#9 du dataset ($y_9 = 44{,}0$ k\texteuro).
\end{enumerate}

% ──────────────────────────────────────────────────────────────────────────
\exercice{\exointermediaire\quad --- Choisir le meilleur mod\`{e}le}

Deux mod\`{e}les concurrents pr\'{e}disent le salaire \`{a} partir de l'exp\'{e}rience seule :

\begin{center}
\renewcommand{\arraystretch}{1.1}
\begin{tabular}{l l}
    \toprule
    \textbf{Mod\`{e}le A :} & $\yhat = 2{,}0 \times \texttt{experience} + 32$ \\
    \textbf{Mod\`{e}le B :} & $\yhat = 3{,}0 \times \texttt{experience} + 25$ \\
    \bottomrule
\end{tabular}
\end{center}

Utilisez les 6 premiers d\'{e}veloppeurs du dataset (voir tableau en haut de page).

\begin{enumerate}[label=\alph*)]
    \item Calculez les pr\'{e}dictions $\yhat_i$ du mod\`{e}le A pour chaque d\'{e}veloppeur.

    \item Calculez les pr\'{e}dictions $\yhat_i$ du mod\`{e}le B pour chaque d\'{e}veloppeur.

    \item Calculez le MSE du mod\`{e}le A et le MSE du mod\`{e}le B.

    \item Quel mod\`{e}le est le meilleur selon le MSE ? Le r\'{e}sultat vous surprend-il ?
\end{enumerate}

% ──────────────────────────────────────────────────────────────────────────
\exercice{\exointermediaire\quad --- Construire la matrice \`{a} partir de donn\'{e}es textuelles}

Voici les profils de 4 d\'{e}veloppeurs :

\begin{itemize}
    \item \textbf{Alice} : 3 ans d'exp\'{e}rience, 2 langages, Bac+3, PME (150 employ\'{e}s),
        80\% remote, salaire 34 k\texteuro.
    \item \textbf{Bob} : 7 ans, 5 langages, Bac+5, scale-up (800 employ\'{e}s),
        50\% remote, salaire 50 k\texteuro.
    \item \textbf{Clara} : 0 an, 1 langage, Bac, micro-entreprise (20 employ\'{e}s),
        100\% remote, salaire 24 k\texteuro.
    \item \textbf{David} : 12 ans, 6 langages, Bac+5, grand groupe (3000 employ\'{e}s),
        30\% remote, salaire 58 k\texteuro.
\end{itemize}

\begin{enumerate}[label=\alph*)]
    \item Construisez la matrice de features $\mX \in \mathbb{R}^{? \times ?}$.
        Pr\'{e}cisez les dimensions.

    \item Construisez le vecteur cible $\vy \in \mathbb{R}^{?}$.

    \item Que vaut l'\'{e}l\'{e}ment $x_{2,4}$ ? Que repr\'{e}sente-t-il concr\`{e}tement ?

    \item En utilisant $\vw = \begin{pmatrix} 2 \\ 1.5 \\ 3 \\ 0.005 \\ 0.02 \end{pmatrix}$ et $b = 15$,
        calculez la pr\'{e}diction pour Bob ($\yhat_2 = \vw^T \vx_2 + b$).
\end{enumerate}

% ══════════════════════════════════════════════════════════════════════════
\section*{Exercices Python guid\'{e}s --- Niveau interm\'{e}diaire}
% ══════════════════════════════════════════════════════════════════════════

% ──────────────────────────────────────────────────────────────────────────
\exercice{\exointermediaire\quad --- Cr\'{e}er et manipuler des vecteurs NumPy \hfill\normalfont\small\textsf{[Python guid\'{e}]}}

Compl\'{e}tez le code suivant dans Google Colab :

\begin{pythonbox}[Squelette --- Vecteurs NumPy]
\begin{lstlisting}
import numpy as np

# 1. Creer le vecteur du developpeur #3 du dataset
x3 = np.array([______])  # A COMPLETER

# 2. Afficher la dimension
print(f"Dimension de x3 : {______}")  # A COMPLETER

# 3. Acceder a la 3e composante (niveau d'etudes)
niveau = ______  # A COMPLETER
print(f"Niveau d'etudes : {niveau}")

# 4. Calculer la norme du vecteur
# Rappel : ||x|| = sqrt(x1**2 + x2**2 + ... + xp**2)
norme = np.sqrt(np.______(x3 ** 2))  # A COMPLETER
print(f"Norme de x3 : {norme:.2f}")

# 5. Creer le vecteur du developpeur #5
x5 = np.array([______])  # A COMPLETER

# 6. Calculer le produit scalaire de x3 et x5
ps = ______  # A COMPLETER (utiliser @)
print(f"Produit scalaire x3 . x5 = {ps}")

# 7. Calculer la distance entre les deux developpeurs
# distance = ||x3 - x5||
distance = np.sqrt(np.sum((______ - ______) ** 2))
print(f"Distance : {distance:.2f}")
\end{lstlisting}
\end{pythonbox}

% ──────────────────────────────────────────────────────────────────────────
\exercice{\exointermediaire\quad --- Construire la matrice du dataset \hfill\normalfont\small\textsf{[Python guid\'{e}]}}

\begin{pythonbox}[Squelette --- Matrice NumPy]
\begin{lstlisting}
import numpy as np

# 1. Creer la matrice des 6 premiers developpeurs
#    (5 features seulement, pas le salaire)
X = np.array([
    [0, 2, 3, 150, 80],
    [1, 1, 2, 50, 100],
    ______,  # A COMPLETER : dev #2
    ______,  # A COMPLETER : dev #3
    ______,  # A COMPLETER : dev #4
    ______,  # A COMPLETER : dev #5
])

# 2. Verifier les dimensions
print(f"Shape de X : {______}")  # doit afficher (6, 5)

# 3. Extraire la colonne "experience" (colonne 0)
experience = X[:, ______]  # A COMPLETER
print(f"Experiences : {experience}")

# 4. Extraire la ligne du developpeur #4
dev4 = X[______, :]  # A COMPLETER
print(f"Developpeur #4 : {dev4}")

# 5. Creer le vecteur des salaires
y = np.array([28.0, 26.5, ______, ______, ______, ______])

# 6. Calculer la transposee et afficher son shape
print(f"Shape de X.T : {______}")  # doit afficher (5, 6)

# 7. Quelle est la valeur de X[2, 3] ?
#    Que represente-t-elle ?
print(f"X[2,3] = {X[2, 3]}")  # interpreter !
\end{lstlisting}
\end{pythonbox}

% ──────────────────────────────────────────────────────────────────────────
\exercice{\exointermediaire\quad --- R\'{e}gression lin\'{e}aire from scratch \hfill\normalfont\small\textsf{[Python guid\'{e}]}}

\begin{pythonbox}[Squelette --- R\'{e}gression lin\'{e}aire]
\begin{lstlisting}
import numpy as np
import matplotlib.pyplot as plt

# Donnees completes (experience, salaire)
experience = np.array([0,1,1,2,2,0,3,1,4,5,5,6,6,7,7,4,
                        8,5,9,10,10,11,12,9,13,11,15,17,20,18])
salaire = np.array([28.0,26.5,35.0,32.0,38.0,25.0,34.5,42.0,
    38.5,44.0,41.0,48.0,43.5,50.0,39.0,47.0,51.0,40.0,
    52.0,55.0,48.5,58.0,60.0,42.0,62.0,54.0,65.0,68.0,72.0,58.0])

# 1. Calculer les moyennes
x_mean = ______  # A COMPLETER
y_mean = ______  # A COMPLETER

# 2. Calculer la pente w (formule des moindres carres)
# w = sum((xi - x_mean) * (yi - y_mean)) / sum((xi - x_mean)**2)
numerateur = np.sum(______)    # A COMPLETER
denominateur = np.sum(______)  # A COMPLETER
w = numerateur / denominateur
print(f"Pente w = {w:.4f}")

# 3. Calculer le biais b
b = ______  # A COMPLETER (formule : y_mean - w * x_mean)
print(f"Biais b = {b:.4f}")

# 4. Calculer les predictions pour tous les developpeurs
y_pred = ______  # A COMPLETER (formule : w * experience + b)

# 5. Calculer le MSE
mse = np.mean(______)  # A COMPLETER
print(f"MSE = {mse:.2f}")

# 6. Tracer le graphique
fig, ax = plt.subplots(figsize=(8, 5))
ax.scatter(experience, salaire, c='blue', label='Donnees')
x_line = np.linspace(0, 21, 100)
ax.plot(x_line, ______, 'r--', label=f'y={w:.2f}x+{b:.2f}')
ax.set_xlabel("Experience (annees)")
ax.set_ylabel("Salaire (kEUR)")
ax.set_title("Regression lineaire from scratch")
ax.legend()
ax.grid(True, alpha=0.3)
plt.show()
\end{lstlisting}
\end{pythonbox}

% ══════════════════════════════════════════════════════════════════════════
\section*{Exercices Python autonomes --- Niveau avanc\'{e}}
% ══════════════════════════════════════════════════════════════════════════

% ──────────────────────────────────────────────────────────────────────────
\exercice{\exoavance\quad --- Explorer le dataset salaire \hfill\normalfont\small\textsf{[Python autonome]}}

\'{E}crivez un programme Python complet (sans squelette) qui r\'{e}alise les t\^{a}ches suivantes
\`{a} partir du dataset fil rouge (30 d\'{e}veloppeurs, 6 colonnes) :

\begin{enumerate}[label=\alph*)]
    \item Chargez le dataset complet dans un tableau NumPy.
        Affichez ses dimensions et les 5 premi\`{e}res lignes.

    \item Calculez et affichez pour chaque colonne :
        la \textbf{moyenne}, l'\textbf{\'{e}cart-type}, le \textbf{minimum} et le \textbf{maximum}.

    \item Calculez la \textbf{corr\'{e}lation} entre chaque feature et le salaire.
        \emph{Rappel :} $r_{xy} = \frac{\sum (x_i - \bar{x})(y_i - \bar{y})}
        {\sqrt{\sum (x_i - \bar{x})^2 \cdot \sum (y_i - \bar{y})^2}}$
        \quad (ou utilisez \texttt{np.corrcoef}).

    \item Tracez un graphique avec 5 sous-figures (une par feature) montrant
        chaque feature en abscisse et le salaire en ordonn\'{e}e.
        Quelle feature semble la plus corr\'{e}l\'{e}e au salaire visuellement ?
\end{enumerate}

\begin{attentionbox}[Consignes]
    \begin{itemize}
        \item Utilisez \textbf{uniquement} NumPy et Matplotlib (pas de pandas).
        \item Commentez votre code en fran\c{c}ais.
        \item Affichez des r\'{e}sultats lisibles avec \texttt{f-strings} format\'{e}es.
    \end{itemize}
\end{attentionbox}

% ──────────────────────────────────────────────────────────────────────────
\exercice{\exoavance\quad --- Tester plusieurs mod\`{e}les \hfill\normalfont\small\textsf{[Python autonome]}}

Vous \^{e}tes data scientist junior. Votre manager vous demande de comparer
3 mod\`{e}les lin\'{e}aires simples (exp\'{e}rience $\to$ salaire) :

\begin{center}
\renewcommand{\arraystretch}{1.2}
\begin{tabular}{c l c c}
    \toprule
    \textbf{Mod\`{e}le} & \textbf{Formule} & $w_1$ & $w_0$ \\
    \midrule
    A & $\yhat = 1{,}5 \times x + 33$ & 1.5 & 33 \\
    B & $\yhat = 2{,}07 \times x + 31{,}27$ & 2.07 & 31.27 \\
    C & $\yhat = 2{,}8 \times x + 28$ & 2.8 & 28 \\
    \bottomrule
\end{tabular}
\end{center}

\'{E}crivez un programme Python complet qui :

\begin{enumerate}[label=\alph*)]
    \item Charge le dataset et extrait les colonnes \texttt{experience} et \texttt{salaire}.

    \item Pour chacun des 3 mod\`{e}les, calcule les pr\'{e}dictions $\hat{\vy}$
        et le MSE associ\'{e}. Affichez un tableau r\'{e}capitulatif.

    \item Trace un \textbf{unique graphique} montrant :
        \begin{itemize}
            \item les 30 points du dataset (nuage de points),
            \item les 3 droites de r\'{e}gression (avec une couleur et un label diff\'{e}rents),
            \item une l\'{e}gende indiquant le MSE de chaque mod\`{e}le.
        \end{itemize}

    \item R\'{e}pondez (en commentaire ou \texttt{print}) :
        quel mod\`{e}le est le meilleur ? Pourquoi \'{e}tait-ce pr\'{e}visible ?
\end{enumerate}

\begin{intuitionbox}[Indice]
    Le mod\`{e}le B correspond \`{a} la droite des moindres carr\'{e}s calcul\'{e}e en cours.
    Est-ce forc\'{e}ment le meilleur ?
\end{intuitionbox}

% ══════════════════════════════════════════════════════════════════════════
\vfill
\separateur

\begin{center}
    \small\sffamily
    \textbf{R\'{e}partition :}\quad
    4 exercices \starfull{ipssiVert} fondamentaux \quad
    3 exercices \starfull{ipssiOrange}\starfull{ipssiOrange} interm\'{e}diaires \quad
    3 guid\'{e}s Python \starfull{ipssiOrange}\starfull{ipssiOrange} \quad
    2 autonomes \starfull{ipssiRouge}\starfull{ipssiRouge}\starfull{ipssiRouge}
\end{center}

\begin{center}
    {\footnotesize\sffamily\textcolor{gray}{%
        \textcopyright\ 2025 Mohamed EL AFRIT --- IPSSI \enspace\textbullet\enspace
        \url{www.mohamedelafrit.com} \enspace\textbullet\enspace
        CC BY-NC-SA 4.0
    }}
\end{center}

\end{document}
